% Report template for Project 1 — Factor Z
\documentclass[11pt,a4paper]{article}
\usepackage[utf8]{inputenc}
\usepackage[T1]{fontenc}
\usepackage{graphicx}
\usepackage{booktabs}
\usepackage{siunitx}
\usepackage{longtable}
\usepackage{caption}
\usepackage{hyperref}
\usepackage{geometry}
\geometry{margin=1in}

\title{Projecto 1 — Cálculo do Factor de Compressibilidade do Gás Natural (Z)}
\author{Grupo — Engenharia de Reservatórios I \\\small{Docente: Geraldo Ramos}}
\date{\today}

\begin{document}
\maketitle
\begin{abstract}
Resumo breve do projecto e dos objectivos.
\end{abstract}

\section{Introdução}
Contexto e objectivos.

\section{Metodologia}
Descrição das correlações implementadas:
\begin{itemize}
  \item Gás Ideal (Z = 1)
  \item Hall--Yarborough (Newton--Raphson)
  \item Dranchuk \\& Abou--Kassem (iteração em densidade reduzida)
  \item Correções pseudo-críticas: Standing, Sutton
  \item Correcções de contaminantes: Wichert--Aziz; Carr--Kobayashi--Burrows
\end{itemize}

\section{Resultados}
Importe a tabela CSV gerada pelo programa (ex.: `results.csv`) e inclua-a aqui. Exemplo com tabular:

\begin{table}[ht]
\centering
\caption{Exemplo de resultados — pressão vs Z}
\begin{tabular}{r r r}
\toprule
Pressão (psia) & Z (Hall--Yarborough) & Z (DAK) \\
\midrule
1000 & 0.8723 & 0.8731 \\
2000 & 0.8902 & 0.8908 \\
3000 & 0.8510 & 0.8509 \\
\bottomrule
\end{tabular}
\end{table}

\section{Propriedades do gás}
Inclua viscosidade (Lee--González--Eakin / Lucas), Bg e Eg, com unidades escolhidas.

\section{Discussão e Conclusões}
Discussão dos resultados e comparação entre métodos.

\section{Instruções para replicação}
O código está em `src/`. Para executar a GUI (refatorada):
\begin{verbatim}
python -m src.main
\end{verbatim}
Gere uma tabela CSV via a função de exportação e inclua-a no relatório.

\end{document}
